% Generated draft: requires mathematical and source review.
\documentclass[11pt]{article}
\usepackage[margin=25mm]{geometry}
\usepackage{fontspec}
\setmainfont{Latin Modern Roman}
\usepackage{amsmath,amssymb,mathtools,unicode-math}
\setmathfont{Latin Modern Math}
\usepackage{enumitem,needspace,xurl,hyperref}
\setlength{\parindent}{0pt}
\setlength{\parskip}{5pt}
\setlength{\emergencystretch}{3em}
\newcommand{\N}{\mathbb{N}}
\newcommand{\Z}{\mathbb{Z}}
\newcommand{\Q}{\mathbb{Q}}
\newcommand{\R}{\mathbb{R}}
\newcommand{\C}{\mathbb{C}}
\newcommand{\F}{\mathbb{F}}
\newcommand{\A}{\mathbb{A}}
\newcommand{\PP}{\mathbb{P}}
\newcommand{\E}{\mathbb{E}}
\newcommand{\eps}{\varepsilon}
\newcommand{\dd}{\,\mathrm{d}}
\begin{document}
\section*{UnsolvedMath ID: 20000832}
\textbf{Source problem code:} catalogue entry
\subsection{Definitions and mathematical statement}
The source considers the Hilbert function $(1,4,10,a)$, $6\leq a\leq10$, says the general Artinian algebra with that function is nonsmoothable, and asks for the generic point of the component containing it. Smoothability means lying in the closure of the locus of reduced points in the relevant Hilbert scheme. The statement does not distinguish the generic algebra over a residue field from the scheme structure at a generic point; the supplied aid makes that distinction in characteristic zero.

Let $k$ have characteristic zero, set $S=k[y_1,y_2,y_3,y_4]$, $\mathfrak m=(y_1,\ldots,y_4)$ and $\dim_k S_3=20$. For a subspace $W\subset S_3$ of dimension $20-a$, put
\[
 I_W=(W)+\mathfrak m^4,\qquad A_W=S/I_W.
\]
This has length $15+a$ and Hilbert function $(1,4,10,a)$. Allowing its support to translate in affine four-space gives the very-compressed locus
\[
 B_a=\A^4_k\times\operatorname{Gr}(20-a,S_3)
 \cong\A^4_k\times\operatorname{Gr}(a,20),
\]
where the latter notation uses $a$-dimensional quotients.

For $a=6,7,8,9$, the aid attributes to Jelisiejew's 2024 theorem the statement that $B_a$ is the reduced support of a generically nonreduced component of $\operatorname{Hilb}^{15+a}(\A^4_k)$. If $\eta$ is its topological generic point and $K_a=k(B_a)$, its generic algebra, in coordinates centered at its support, is
\[
 K_a[y_1,\ldots,y_4]/\bigl((W_\eta)+(y_1,\ldots,y_4)^4\bigr),
\]
with $W_\eta$ the generic $(20-a)$-plane in the cubic space. The Hilbert-scheme local ring at $\eta$ is nevertheless a nonreduced Artinian ring with residue field $K_a$, not just that field. The record does not identify its nilpotent structure or length.

For $a=10$, the same quotient describes the generic algebra of the very-compressed locus, whose dimension is $4+10(20-10)=104$. The aid leaves open whether this locus itself is a component; neither its parameterization nor nonsmoothability alone resolves that component question.

Uniformly for $6\leq a\leq10$, the aid reports that the generic quotient is level of type $a$: its socle is exactly its degree-three part, of dimension $a$. The locus where a socle element occurs below degree three has codimension at least $4a-9$ in the Grassmannian parameter space. At an origin-supported point, with the usual degree convention for graded homomorphisms, it also states
\[
 \operatorname{Hom}_S(I_W,A_W)_{>0}=0,\qquad
 \operatorname{Hom}_S(I_W,A_W)_0
 \cong\operatorname{Hom}_k(W,S_3/W).
\]
Translations inject as four degree-minus-one tangent directions. Consequently all infinitesimal directions normal to the reduced very-compressed locus have negative weight. This tangent-weight information does not by itself give the generic nonreduced local ring or settle the $a=10$ component issue.
\subsection{Short English statement}
For four of the five values, the supplied theorem identifies the component's reduced support and generic quotient, while its generic local ring still has nilpotents. At $a=10$, the same family has a clear generic algebra but is not established here to be a component. The aid also gives levelness and tangent-weight constraints without determining the nilpotent structure.
\subsection{Sources}
\begin{itemize}
\item Local UnsolvedMath registry, permanent ID $20000832$, source code \texttt{catalogue entry}, supplied complete source record and research summary.
\item American Institute of Mathematics, \emph{Components of Hilbert Schemes}, Problem 31, supplied URL \url{https://aimath.org/WWN/hilbertschemes/hilbertschemes.pdf}; canonical local record \texttt{aim-arithmetic-geometry-notes.json}, \texttt{notes[77]}. No PDF was retrieved.
\item Jelisiejew's 2024 theorem, as named in the supplied summary for $a=6,7,8,9$. No title, identifier, theorem number or full bibliographic entry is supplied. The original literature field is empty.
\item Research attempt 1: partial result; problem partially solved at the recorded run; supplied rigor label ``proved'', novelty confidence ``low''. The machine-generated aid is not peer reviewed and requires independent verification before citation; the remaining-open assessment for $a=10$ was not independently updated.
\end{itemize}
\end{document}
